% !PRINTED 523
\weberpage{523}

%<<CONT
\gesperrt{cients are integers in $\Omega$, and passes over into $\chi$ through multiplication with a rational functional which is a unit.}

\begin{center}
\S.\,143.

\gesperrt{The prime factors of the numbers of the field $\Omega$.}
\end{center}

Since the numbers are also contained among the functionals of the field $\Omega$, it follows from the results of \S.\,141 that the integers of the field $\Omega$ can likewise be decomposed into prime factors, but into prime factors which in general are not numbers but functionals. It is, however, not excluded that in special cases numbers too may occur among the prime functionals, which are then called \gesperrt{primes of the field $\Omega$}.

All equations between functionals are in the last instance identical equations between integral rational functions, and resolve themselves into a sequence of equations between numbers. They therefore remain correct when other symbols are put in place of the variables.

From this it follows that an integral functional, a unit, a prime functional do not cease to be integral functionals, units, prime functionals when other symbols are put in place of the variables.

If, then, we decompose a number into its prime factors, and if functionals occur among these prime factors, the variables in the equation representing this decomposition can be replaced by arbitrary other variables, and from this follows the generalisation:

1. \gesperrt{The prime factors of an integral functional $\omega$ can be so represented that they do not contain the variables on which $\omega$ depends.}

For every integral functional $\omega$ is a divisor of numbers, even of rational numbers, for example of the absolute norm of $\omega$. The prime factors of $\omega$ are therefore to be sought among the prime factors of one of these numbers, and can accordingly be represented by variables which are different from those occurring in $\omega$.
